%----------------------------------------------------
\section{Theoretical Foundations of \textsc{Voila}}
\label{sec:theory}
%----------------------------------------------------

We formalize adaptive fidelity selection as a \textbf{cost-sensitive Bayesian decision problem},
show that \textsc{Voila} implements a \emph{near-Bayes-optimal} sequential greedy policy when its
probability estimates are calibrated, and quantify the resulting utility gap on unseen
(test) data.

\vspace{-0.3em}
\subsection{Bayes-Optimal Fidelity Selector}
\label{sec:bayes_optimal}

Let $\mathcal{F}=\{f_1,\dots,f_K\}$ denote the
available visual fidelities
(e.g., caption, JPEG $q{=}10$, full\_res),
ordered by increasing cost $c_1 < c_2 < \cdots < c_K$, where each $c_k > 0$ represents
latency, bandwidth, or storage retrieval cost.

For a question~$q$, let
%
\[
p_k(q)\;=\;\Pr\!\bigl[\,\text{answer correct} \mid q, f_k\bigr]
\]
%
denote the \emph{true} probability of answering correctly using fidelity~$f_k$.
Define the \textbf{utility} of selecting fidelity~$f_k$ as
\[
U_k(q)\;=\;p_k(q)\;-\;\lambda c_k,
\]
where $\lambda > 0$ is a user-chosen trade-off coefficient that controls the relative importance
of accuracy versus cost.

The \textbf{Bayes-optimal} decision rule is
\begin{equation}
f^\star(q)\;=\;\argmax_{k \in \{1,\dots,K\}} U_k(q).
\tag{1}
\label{eq:bayes_optimal}
\end{equation}

This rule assumes \emph{oracle access} to the true probabilities $p_k(q)$ for all fidelities.
In practice, these probabilities are unknown and must be estimated from data.

\vspace{-0.3em}
\subsection{Learning Utility via Question-Conditioned Calibration}
\label{sec:calibration}

Because $p_k(q)$ is unknown at test time,
\textsc{Voila} learns an estimator via a two-stage pipeline:

\paragraph{Stage 1: Gradient-Boosted Regression.}
For each fidelity~$f_k$, we extract question features $\phi(q) \in \mathbb{R}^d$
(TF-IDF embeddings, question length, WH-word indicators, etc.)
and train a Gradient Boosting Regressor (GBR) to predict a raw score:
\[
\tilde z_k(q) \;=\; \text{GBR}_k\bigl(\phi(q)\bigr).
\]
The GBR is trained on labeled data $\{(q_i, y_{i,k})\}$, where $y_{i,k} \in \{0,1\}$ indicates
whether the model answered question~$q_i$ correctly using fidelity~$f_k$.

\paragraph{Stage 2: Isotonic Calibration.}
The raw scores $\tilde z_k$ are typically poorly calibrated. We apply
\emph{isotonic regression}~\cite{nuesch1991order}
to map them to calibrated probabilities:
\[
\hat p_k(q)\;=\;
g_k\!\bigl(\sigma\!\bigl(\tilde z_k(q)\bigr)\bigr),
\]
where $\sigma$ is the logistic sigmoid and $g_k: [0,1] \to [0,1]$ is a monotone non-decreasing
function fitted on a held-out calibration set to minimize calibration error
(e.g., expected calibration error or Brier score).

\paragraph{Deployment Selector.}
At test time, \textsc{Voila} computes the estimated utility for each fidelity and selects
\begin{equation}
\hat f(q)
\;=\;
\argmax_{k \in \{1,\dots,K\}}
\bigl\{\,
\hat p_k(q)\;-\;\lambda c_k
\bigr\}.
\tag{2}
\label{eq:voila_selector}
\end{equation}

Crucially, this decision is made \emph{before retrieving any visual fidelity}, using only
question features $\phi(q)$. This distinguishes \textsc{Voila} from post-retrieval model
routing approaches, which assume all context is already available.

\vspace{-0.3em}
\subsection{Optimality Under Perfect Calibration}
\label{sec:optimality}

\begin{proposition}[Exact Optimality]
If $\hat p_k(q)=p_k(q)$ for all $k \in \{1,\dots,K\}$ and all queries $q$,
then the selector~\eqref{eq:voila_selector} coincides with the Bayes-optimal rule~\eqref{eq:bayes_optimal}.
\end{proposition}

\begin{proof}
Identical utility functions yield identical maximizers.
\end{proof}

\noindent
This result establishes that calibration quality directly determines decision optimality.

\vspace{-0.3em}
\subsection{Regret with Imperfect Calibration}
\label{sec:regret}

In practice, calibration is imperfect. We quantify the decision quality degradation
via a \emph{regret bound}.

\begin{definition}[Uniform Calibration Error]

\[
\varepsilon\;=\;
\sup_{q \in \mathcal{Q}}\max_{k \in \{1,\dots,K\}}\bigl|\,\hat p_k(q)-p_k(q)\bigr|,
\]
where $\mathcal{Q}$ is the distribution of test queries.
\end{definition}

\begin{theorem}[Utility Regret Bound]
\label{thm:regret}
Let $U^\star(q)=\max_{k}U_k(q)$
and $\hat U(q)=U_{\hat f(q)}(q)$ denote the optimal and achieved utilities, respectively.
For every question $q$,
\[
U^\star(q)\;-\;\hat U(q)\;\le\;2\varepsilon.
\]
Hence the population-level expected utility gap satisfies
$\mathbb{E}_{q \sim \mathcal{Q}}[U^\star(q)-\hat U(q)]\le 2\varepsilon.$
\end{theorem}

\begin{proof}
Let $k^\star=f^\star(q)$ be the Bayes-optimal fidelity and $\hat k = \hat f(q)$ be the fidelity
selected by \textsc{Voila}.
By definition of $\argmax$,
\begin{equation}
\hat p_{\hat k}(q) - \lambda c_{\hat k}
\;\ge\;
\hat p_{k^\star}(q) - \lambda c_{k^\star}.
\label{eq:argmax_ineq}
\end{equation}
The true utility gap is
\begin{align*}
U^\star(q) - \hat U(q)
&= \bigl(p_{k^\star}(q) - \lambda c_{k^\star}\bigr)
- \bigl(p_{\hat k}(q) - \lambda c_{\hat k}\bigr) \\
&= \bigl(p_{k^\star}(q) - p_{\hat k}(q)\bigr)
- \lambda\bigl(c_{k^\star} - c_{\hat k}\bigr).
\end{align*}
Using $|p_k(q) - \hat p_k(q)| \le \varepsilon$ for all $k$, we have
\[
p_{k^\star}(q) \le \hat p_{k^\star}(q) + \varepsilon
\quad\text{and}\quad
p_{\hat k}(q) \ge \hat p_{\hat k}(q) - \varepsilon.
\]
Substituting,
\begin{align*}
U^\star(q) - \hat U(q)
&\le \bigl(\hat p_{k^\star}(q) + \varepsilon - (\hat p_{\hat k}(q) - \varepsilon)\bigr)
- \lambda\bigl(c_{k^\star} - c_{\hat k}\bigr) \\
&= \bigl(\hat p_{k^\star}(q) - \lambda c_{k^\star}\bigr)
- \bigl(\hat p_{\hat k}(q) - \lambda c_{\hat k}\bigr) + 2\varepsilon \\
&\le 2\varepsilon,
\end{align*}
where the last inequality follows from~\eqref{eq:argmax_ineq}.
\end{proof}

\paragraph{Interpretation.}
The bound is tight and linear in the calibration error: every additional unit of miscalibration
can degrade utility by at most two units. This provides a direct, actionable link between
learning quality (calibration) and system-level performance (decision utility).

\vspace{-0.4em}
\subsection{Generalization to Unseen Questions}
\label{sec:generalization}

Isotonic regression enjoys well-known uniform convergence guarantees~\cite{yang2018contractionuniformconvergenceisotonic, lim2025review}.
%We adapt results from~\citet{yang2018contractionuniformconvergenceisotonic} to our setting.

\begin{lemma}[Empirical $\!\to\!$ Population Calibration]
\label{lem:gen}
Assume the calibration set $\{(q_i, y_{i,k})\}_{i=1}^N$ is drawn i.i.d.\ from the test
distribution $\mathcal{Q}$, and that the raw scores $\tilde z_k(q)$ are bounded in $[a,b]$.
With probability at least $1-\delta$,
\[
\varepsilon
\;\le\;
\hat\varepsilon
\;+\;
\underbrace{O\!\left(\sqrt{\tfrac{\log(K/\delta)}{N}}\right)}_{\text{statistical gap}},
\]
where $\hat\varepsilon$ is the empirical calibration error on the held-out calibration set,
and $K$ is the number of fidelities.
\end{lemma}

\begin{proof}[Proof Sketch]
Isotonic regression minimizes the squared loss over monotone functions, which forms a
class with finite VC dimension. Standard uniform convergence results for empirical risk
minimization~\citep{vapnik1968uniform} yield the stated rate. The $\log K$ term accounts
for union bound over $K$ fidelity-specific calibrators.
\end{proof}

\noindent
Combining Lemma~\ref{lem:gen} with Theorem~\ref{thm:regret} yields a
\textbf{distribution-free test-time regret bound} that decays as $O(N^{-1/2})$,
verified empirically in Section~\ref{sec:results}.

\vspace{-0.3em}
\subsection{Sequential Greedy Policy}
\label{sec:greedy_policy}

The theory in Sections~\ref{sec:bayes_optimal}--\ref{sec:generalization} analyzes
\emph{single-step} fidelity selection: choose one fidelity and commit.
In practice, \textsc{Voila} implements a \emph{sequential greedy escalation} policy
(Algorithm~1 in the main text): start with the cheapest fidelity $f_1$, and iteratively
escalate to $f_{k+1}$ if the incremental value of information
\[
\text{VOI}(f_{k+1} \mid f_k)
\;=\;
\bigl(\hat p_{k+1}(q) - \hat p_k(q)\bigr) - \lambda c_{k+1}
\]
exceeds a threshold $\tau \ge 0$.

This greedy policy is \emph{myopic}---it does not plan ahead for future escalations---but
enjoys the following guarantee:

\begin{corollary}[Regret of Greedy Sequential Policy]
\label{cor:greedy_regret}
Assume (i) calibrated probabilities are non-decreasing across fidelity levels,
$\hat p_{f_1}(q) \le \hat p_{f_2}(q) \le \cdots \le \hat p_{f_K}(q)$
(enforced by the cross-fidelity isotonic pass in \S\ref{sec:method}),
and (ii) marginal accuracy gains are non-increasing (diminishing returns):
$\hat p_{f_{k+1}}(q) - \hat p_{f_k}(q) \ge \hat p_{f_{k+2}}(q) - \hat p_{f_{k+1}}(q)$
for all $k$.
Let $\hat f_{\text{greedy}}(q)$ denote the fidelity selected by the legacy
sequential escalation policy (\S\ref{sec:method})
with $\tau = 0$. Then $\hat f_{\text{greedy}}(q) = \hat f(q)$ (the argmax selector of
Eq.~\eqref{eq:voila_selector}), and consequently
\[
U^\star(q) - U_{\hat f_{\text{greedy}}(q)}(q)
\;\le\;
2\varepsilon.
\]
\end{corollary}

\begin{proof}
We show that $\mathrm{VOI}(f_{k+1} \mid f_k) \le 0$ implies
$\mathrm{VOI}(f_{k+2} \mid f_{k+1}) \le 0$, so the utility function
$U_k(q) = \hat p_{f_k}(q) - \lambda c_k$ is unimodal and the greedy break
is never premature.

By assumption (ii), $\delta_{k+1} := \hat p_{f_{k+1}}(q) - \hat p_{f_k}(q) \ge
\delta_{k+2} := \hat p_{f_{k+2}}(q) - \hat p_{f_{k+1}}(q)$.
Suppose $\mathrm{VOI}(f_{k+1} \mid f_k) = \delta_{k+1} - \lambda c_{k+1} \le 0$,
i.e., $\delta_{k+1} \le \lambda c_{k+1}$.
By costs strictly increasing and diminishing returns:
\[
\mathrm{VOI}(f_{k+2} \mid f_{k+1})
= \delta_{k+2} - \lambda c_{k+2}
\le \delta_{k+1} - \lambda c_{k+1}
\le 0,
\]
where the first inequality uses $\delta_{k+2} \le \delta_{k+1}$ and $c_{k+2} > c_{k+1}$.
Hence the greedy terminates at the unique peak of $U_k(q)$,
which coincides with $\hat f(q)$ from Eq.~\eqref{eq:voila_selector},
and Theorem~\ref{thm:regret} applies.
\end{proof}

\noindent
\textbf{Remark.}
When calibrated probabilities are non-monotone across fidelities for a given query
(possible if the cross-fidelity isotonic pass is omitted), the greedy break may
terminate at a local maximum below the global optimum.
In such cases one should fall back to the exhaustive argmax of Eq.~\eqref{eq:voila_selector},
which always satisfies Theorem~\ref{thm:regret} without additional assumptions.

\noindent
In practice, setting $\tau > 0$ trades off between exploration (fewer escalations, lower cost)
and exploitation (higher accuracy). The optimal threshold depends on the cost regime $\lambda$
and is chosen via cross-validation on a held-out set.

\vspace{-0.3em}
\subsection{Implications for Adaptive Fidelity Selection}
\label{sec:implications}

\textbf{(i) Fixed fidelity is provably suboptimal.}
For a heterogeneous query mix, any single $f_k$ incurs
additive cost $\lambda\Delta c$ on ``easy'' questions (where lower fidelity suffices)
or accuracy loss on ``hard'' ones (where higher fidelity is required),
whereas \textsc{Voila} closes the gap up to $2\varepsilon$.

\textbf{(ii) VLM confidence thresholding lacks guarantees.}
Threshold heuristics that rely on raw VLM token probabilities ignore acquisition costs,
are uncalibrated across fidelities (as shown empirically in Figure~3),
and provide no regret bounds. In contrast, \textsc{Voila}'s calibrated question-conditioned
predictors yield interpretable, cost-aware decisions with formal guarantees.

\textbf{(iii) Calibration $\Rightarrow$ efficiency \& accuracy.}
By tethering decision quality to calibration error, \textsc{Voila} turns an ML
problem (probability calibration) into a \emph{system-level performance knob}:
better calibration yields proportionally better cost-accuracy trade-offs.
This provides a principled optimization target for system designers.

\textbf{(iv) Pre-retrieval vs.\ post-retrieval optimization.}
Unlike model routing approaches that select \emph{which model} to apply after
retrieving all context, \textsc{Voila} selects \emph{which context to retrieve}
before any visual data is accessed. These two problems are orthogonal and complementary:
\textsc{Voila} can be combined with model routing to jointly optimize both retrieval
and computation.

\paragraph{Take-away.}
A lightweight two-stage pipeline (GBR + isotonic calibration) gives \textsc{Voila}
\emph{provable} near-optimality guarantees with explicit regret bounds,
providing theoretical backing for the consistent empirical gains observed
over fixed-fidelity and threshold baselines across diverse multimodal VQA benchmarks.

\vspace{-0.2em}
\subsection{Connection to Empirical Observations}
\label{sec:theory_empirics}

Our theoretical framework directly addresses the failure modes identified in Section~2:

\textbf{Observation 1 (Query-dependent information requirements):}
The Bayes-optimal rule~\eqref{eq:bayes_optimal} explicitly accounts for heterogeneous
$p_k(q)$ across queries, selecting higher fidelity only when it increases utility.
Figure~2 shows this heterogeneity empirically; Theorem~\ref{thm:regret} shows
\textsc{Voila} adapts near-optimally.

\textbf{Observation 2 (Unreliable VLM confidence):}
Unlike threshold-based methods that rely on post-retrieval VLM token probabilities
(Figure~3), \textsc{Voila} learns \emph{cross-fidelity performance predictors} from
question features alone, bypassing the VLM's miscalibration entirely.

\textbf{Observation 3 (Model scaling cannot compensate for missing information):}
Figure~4 shows that larger models do not overcome low-fidelity bottlenecks.
Our framework formalizes this: $p_k(q)$ is bounded by the \emph{information content}
of fidelity~$f_k$, independent of model capacity. Adaptive fidelity selection addresses
this fundamental limit, whereas model routing does not.

\textbf{Observation 4 (Calibration enables reliable decisions):}
Proposition~\ref{sec:optimality} and Theorem~\ref{sec:regret} formalize the empirical
observation that calibrated predictors enable cost-aware decisions with quantifiable
performance guarantees.